\documentclass[11pt]{article}
\usepackage[letterpaper,margin=1in]{geometry}
\usepackage[T1]{fontenc}
\usepackage{lmodern}
\usepackage{microtype}
\usepackage{amsmath,amssymb,amsthm,mathtools}
\usepackage{booktabs,tabularx,longtable,array}
\usepackage{enumitem}
\usepackage{tocloft}
\setlength{\cftbeforesecskip}{3pt}
\usepackage{xcolor}
\usepackage{hyperref}
\hypersetup{colorlinks=true,linkcolor=blue!45!black,citecolor=blue!45!black,urlcolor=blue!45!black,
 pdftitle={A fluctuation theory for bounded-batch species discovery},pdfauthor={Christopher D. Long}}
\usepackage{fancyhdr}
\pagestyle{fancy}
\fancyhf{}
\fancyhead[L]{\small Bounded-batch species discovery}
\fancyhead[R]{\small Working draft}
\fancyfoot[C]{\thepage}
\setlength{\headheight}{14pt}
\setlength{\emergencystretch}{2em}
\numberwithin{equation}{section}
\newtheorem{theorem}{Theorem}[section]
\newtheorem{lemma}[theorem]{Lemma}
\newtheorem{proposition}[theorem]{Proposition}
\newtheorem{corollary}[theorem]{Corollary}
\theoremstyle{definition}
\newtheorem{definition}[theorem]{Definition}
\newtheorem{example}[theorem]{Example}
\theoremstyle{remark}
\newtheorem{remark}[theorem]{Remark}
\newcommand{\E}{\mathbb E}
\newcommand{\Pp}{\mathbb P}
\newcommand{\N}{\mathbb N}
\newcommand{\R}{\mathbb R}
\newcommand{\ind}{\mathbf 1}
\newcommand{\Var}{\operatorname{Var}}
\newcommand{\Cov}{\operatorname{Cov}}
\newcommand{\Poi}{\operatorname{Poi}}
\newcommand{\Exp}{\operatorname{Exp}}
\newcommand{\Law}{\mathcal L}
\newcommand{\dK}{d_{\mathrm K}}
\newcommand{\Normal}{\mathcal N(0,1)}
\newcommand{\dto}{\xrightarrow{\ d\ }}
\newcommand{\pto}{\xrightarrow{\ \Pp\ }}
\newcommand{\phifun}{\psi}
\DeclareMathOperator{\supp}{supp}
\title{A fluctuation theory for\ bounded-batch species discovery}
\author{Christopher D. Long}
\date{November 12, 2026\\[3pt]\small Revised research draft; proof and reference audit}
\begin{document}
\maketitle
\begin{abstract}
Let independent, identically distributed random subsets of a countable species
space have cardinality at most $B$. We prove a central limit theorem for the
number of distinct species after $n$ observations whenever its variance diverges,
without within-batch independence or connectedness. This answers the vertex-count
question in Janson's countable-edge model and extends it to bounded hyperedges.
A positive weighted-overlap representation of a terminal martingale's predictable
quadratic variation gives variance and exponential-moment bounds for that bracket.
Writing $V$ for the Poissonized count variance and $\eta=B^2/V$, we obtain a
uniform Kolmogorov bound of order $[\eta\log(e+\eta^{-1})]^{1/4}$, as well as
the simpler fourth-moment bound $O(\eta^{1/5})$. We also prove
$\Pp(|X-\E X|\ge x)\le2\exp\{-x^2/[2(V+2Bx)]\}$ and a Gaussian
moderate-deviation principle. Product-space inequalities separately give the
sharp comparison $v_0\le V\le Bv_0$, species-dependent variance refinements,
and fractional Bennett bounds. For fixed laws with $\E|A_1|^2<\infty$, an
absolute expansion $\Var(D_n)=V(n)-n m'(n)^2+o(1)$ enables de-Poissonization.
A Lindeberg condition extends the CLT to unbounded batches; a disjoint-block
example shows that a finite second moment alone is insufficient. Further results
cover growing multiple-coverage thresholds and nonnegative weighted discovery.
Joint occupancy-spectrum limits, functional limits, and relative Gaussian-tail
approximations are not asserted.
\end{abstract}

\setcounter{tocdepth}{1}
{\small\tableofcontents}
\clearpage

\section{Model, main theorem, and relation to previous work}\label{sec:intro}

Let $\mathcal S$ be countable and let $A_1,A_2,\ldots$ be independent random
finite subsets of $\mathcal S$ with a common deterministic law $\mu$. Unless
stated otherwise,
\begin{equation}\label{eq:batchbound}
 |A_k|\le B\quad\text{almost surely},\qquad B\in\N.
\end{equation}
Empty batches are allowed. Define
\[
 D_n=\left|\bigcup_{k=1}^n A_k\right|,\qquad
 X(t)=D_{N(t)},
\]
where $N$ is an independent rate-one Poisson process. The law $\mu$ is fixed in
statements concerning $n\to\infty$ unless a triangular array is explicitly
specified. Set
\begin{align}
 \lambda_i&=\Pp(i\in A_1),&
 \lambda_{ij}&=\Pp(i,j\in A_1),\qquad \lambda_{ii}=\lambda_i,\label{eq:rates}\\
 m(t)&=\E X(t)=\sum_i(1-e^{-t\lambda_i}),&
 b_n&=\E D_n=\sum_i(1-(1-\lambda_i)^n),\label{eq:means}\\
 V(t)&=\Var(X(t)),&s_n^2&=\Var(D_n),\label{eq:vars}\\
 v_0(t)&=\sum_i e^{-t\lambda_i}(1-e^{-t\lambda_i}),&
 f(t)&=m'(t)=\sum_i\lambda_i e^{-t\lambda_i}.\label{eq:flux}
\end{align}
Species with $\lambda_i=0$ are removed. The identities
\begin{equation}\label{eq:rows}
 \sum_i\lambda_i=\E|A_1|\le B,\qquad
 \sum_j\lambda_{ij}=\E[\ind_{\{i\in A_1\}}|A_1|]\le B\lambda_i
\end{equation}
will be used repeatedly. All sums over $(i,j)$ are over \emph{ordered} pairs,
including the diagonal, unless indicated otherwise.

\begin{theorem}[Bounded-batch discovery CLT]\label{thm:main}
For every fixed law satisfying \eqref{eq:batchbound},
\[
 s_n^2\longrightarrow\infty
 \quad\Longleftrightarrow\quad V(n)\longrightarrow\infty
 \quad\Longleftrightarrow\quad v_0(n)\longrightarrow\infty.
\]
Whenever these equivalent conditions hold,
\begin{equation}\label{eq:fixedCLT}
 \frac{D_n-b_n}{s_n}\dto\Normal,
 \qquad s_n^2\sim V(n).
\end{equation}
For Poissonized sampling, along any sequence with $V(t)\to\infty$,
\begin{equation}\label{eq:poisCLT}
 \frac{X(t)-m(t)}{\sqrt{V(t)}}\dto\Normal.
\end{equation}
No connectedness, regular variation, or within-batch independence assumption is
required.
\end{theorem}

The Poissonized result is quantitative and uniform over the sampling law; see
Theorem~\ref{thm:BE}. The proof of the fixed-sample result is completed in
Section~\ref{sec:depois}. The two useful concentration mechanisms are kept
separate: the read-$B$ argument gives the sharp universal variance proxy $Bv_0(t)$,
whereas the terminal-martingale argument gives a Bernstein bound with the actual
variance $V(t)$ as its quadratic coefficient and linear scale $2B$.
The logarithmic refinement of the Berry--Esseen bound is
Theorem~\ref{thm:BElog}; its proof uses stopping, Gaussian completion, and a
smoothing lemma, rather than a moment order tending to infinity.
Section~\ref{sec:unbounded} gives a Lindeberg extension beyond bounded batches
and a counterexample to a finite-second-moment-only CLT.

\subsection{The graph problem and attribution}
Janson~\cite[Problem~6.9]{Janson} asks for asymptotic normality of the number of
vertices in the countable-edge model when its variance diverges. Eriksson's
Theorem~14~\cite{ErikssonGraphs} proves a vertex CLT in the eventually forever
connected regime by coupling the vertex process to independent urns; his
Proposition~16 treats isolated edges. Taking $B=2$ in Theorem~\ref{thm:main}
removes that connectivity restriction for every fixed distribution on countably
many vertex pairs. This statement concerns the countable-vertex model, not
additional diffuse dust mechanisms or arbitrary mixtures over random laws.

The independent exponential edge-clock construction already appears in the
cited graph literature. The discovery-derivative estimate used for
 de-Poissonization is the classical infinite-urn estimate of
Dutko~\cite{Dutko}, reproduced with proof in
Eriksson~\cite[Lemmas~29--33, especially Lemma~30]{ErikssonGraphs}.
The product-space inequalities below are instances of classical generalized
H\"older and read-$k$ methods~\cite{Finner,GLSS}; we make no novelty claim for
that general machinery. The crucial model-specific step is instead the
positive, weighted-overlap representation of the \emph{predictable quadratic
variation}. It is this representation, rather than an approximation by
independent species, that yields the general Gaussian theorem.

Eriksson's generalized unseen-species work already establishes concentration
for multiple-coverage counts~\cite[Lemma~29]{ErikssonSpecies}. The present
variance-sensitive inequalities and fluctuation results should therefore not
be described as the first concentration results for incidence data. Nor do
Gaussian limits for oracle counts by themselves construct calibrated prediction
intervals with unknown parameters.

\section{Product-space inequalities with bounded weighted overlap}\label{sec:overlap}

Let $(\xi_k)$ be independent coordinates. A random variable $F_\alpha$ has
support $S_\alpha$ if it is measurable with respect to
$\sigma(\xi_k:k\in S_\alpha)$. A support need not be finite.

\begin{lemma}[Weighted variance drop]\label{lem:variance-drop}
Suppose $a_\alpha\ge0$ and
\begin{equation}\label{eq:overlapcondition}
 \sup_k\sum_{\alpha:k\in S_\alpha}a_\alpha\le L.
\end{equation}
For a finite family of square-integrable $F_\alpha$,
\begin{equation}\label{eq:weightedvariance}
 \Var\left(\sum_\alpha a_\alpha F_\alpha\right)
 \le L\sum_\alpha a_\alpha\Var(F_\alpha).
\end{equation}
If $0\le F_\alpha\le1$, this is at most
$L\E\sum_\alpha a_\alpha F_\alpha$.
The bounded-summand assertion extends to nonnegative countable sums of finite
expectation. The variance inequality also extends to centered sums converging
in $L^2$ whenever its right-hand side is finite.
\end{lemma}
\begin{proof}
Expose the independent coordinates in order. With
$\mathcal G_k=\sigma(\xi_1,\ldots,\xi_k)$, write
\[
 \Delta_kF_\alpha=\E[F_\alpha\mid\mathcal G_k]
                  -\E[F_\alpha\mid\mathcal G_{k-1}].
\]
Independence gives $\Delta_kF_\alpha=0$ when $k\notin S_\alpha$.
Weighted Cauchy--Schwarz gives
\[
 \left(\sum_\alpha a_\alpha\Delta_kF_\alpha\right)^2
 \le L\sum_\alpha a_\alpha(\Delta_kF_\alpha)^2.
\]
Take expectations and sum over $k$, using orthogonality and $L^2$ martingale
convergence. The last assertion for bounded nonnegative summands follows from
$\Var(F_\alpha)\le\E F_\alpha$. For countably many such terms, finite partial
sums $S_J$ satisfy
$\E S_J^2\le L\E S_J+(\E S_J)^2$. Monotone convergence gives an $L^2$ limit,
and the inequalities pass to that limit. Centered-series extensions follow by
the same inequality applied to tails.
\end{proof}

\begin{lemma}[Fractional product inequality]\label{lem:finner}
For nonnegative $G_\alpha$ with supports $S_\alpha$, and weights
$w_\alpha\ge0$ such that
$\sum_{\alpha:k\in S_\alpha}w_\alpha\le1$ for every $k$,
\begin{equation}\label{eq:finner}
 \E\prod_\alpha G_\alpha^{w_\alpha}
 \le\prod_\alpha(\E G_\alpha)^{w_\alpha}.
\end{equation}
The statement holds for every finite family of nonnegative integrable
$G_\alpha$ on countably many independent coordinates. Terms of weight zero
are omitted, and constant functions may be factored out. No finiteness
assumption on the supports is required.
\end{lemma}
\begin{proof}
For finitely many coordinates, integrate one coordinate at a time. Conditional
on all the others, generalized H\"older with exponents $1/w_\alpha$ for the
functions involving that coordinate gives the assertion for that integration;
insert a constant function if the weights sum to less than one. Repeat for the
remaining coordinates. For a finite family on countably many coordinates, first set
$G_{\alpha,R}=G_\alpha\wedge R$. Let
$G_{\alpha,R,K}=\E[G_{\alpha,R}\mid\xi_1,\ldots,\xi_K]$.
Independence makes this function measurable with respect to the coordinates
in $S_\alpha\cap\{1,\ldots,K\}$. The finite-coordinate inequality applies.
For fixed $R$, bounded martingale convergence and dominated convergence let
$K\to\infty$, on both sides of \eqref{eq:finner}. Finally let $R\to\infty$;
monotone convergence applies to the finite product and to each expectation.
This is the product-space form of Finner's inequality~\cite{Finner}.
\end{proof}

Put $\psi(u)=e^u-1-u$ and
\[
 h(u)=(1+u)\log(1+u)-u,\qquad u>-1.
\]

\begin{proposition}[Exponential weighted-overlap bound]\label{prop:overlapmgf}
Under \eqref{eq:overlapcondition}, suppose $0\le F_\alpha\le1$, and put
$S=\sum_\alpha a_\alpha F_\alpha$ with $\E S<\infty$. Then for every real
$\theta$,
\begin{equation}\label{eq:overlapmgf}
 \log\E e^{\theta(S-\E S)}
 \le\frac{\E S}{L}\psi(L\theta).
\end{equation}
\end{proposition}
\begin{proof}
For a finite sum apply Lemma~\ref{lem:finner} with
$w_\alpha=a_\alpha/L$ and $G_\alpha=e^{L\theta F_\alpha}$. Convexity on
$[0,1]$ and $\log(1+u)\le u$ give
\[
 \log\E e^{\theta S}
 \le\sum_\alpha\frac{a_\alpha}{L}
        \log\E e^{L\theta F_\alpha}
 \le\frac{e^{L\theta}-1}{L}\E S.
\]
Center and pass to countable sums. For positive $\theta$ use monotone
convergence before centering; for negative $\theta$ use bounded convergence.
\end{proof}

\begin{corollary}[Read-$B$ variance and exponential bounds]\label{cor:readB}
Suppose each independent coordinate occurs in the supports of at most $B$
variables $Z_i$. Then, for finite sums or the corresponding convergent limits,
\begin{align}
 \Var\left(\sum_i Z_i\right)&\le B\sum_i\Var(Z_i),\label{eq:readvar}\\
 \log\E e^{\theta\sum_i(Z_i-\E Z_i)}
 &\le\frac1B\sum_i\log\E e^{B\theta(Z_i-\E Z_i)}.\label{eq:readmgf}
\end{align}
If $Z_i$ are indicators, with $v=\sum_i\Var(Z_i)<\infty$, then
\begin{equation}\label{eq:readB-bennett}
 \log\E e^{\theta\sum_i(Z_i-\E Z_i)}
 \le\frac vB\psi(B|\theta|),\qquad
 \Pp\left(\left|\sum_i(Z_i-\E Z_i)\right|\ge x\right)
 \le2\exp\left\{-\frac vB h\left(\frac xv\right)\right\}.
\end{equation}
The zero-variance case is interpreted separately as a constant sum.
\end{corollary}
\begin{proof}
Use Lemmas~\ref{lem:variance-drop} and~\ref{lem:finner} with unit coefficients
and weights $1/B$. For an indicator $Z$, $U=Z-\E Z$ satisfies $|U|\le1$,
$\E U=0$, and $\E|U|^k\le\E U^2$ for $k\ge2$. Expansion of the exponential
therefore yields
$\log\E e^{uU}\le\Var(Z)\psi(|u|)$.
Chernoff optimization for each sign proves \eqref{eq:readB-bennett}.
\end{proof}

\section{Sharp variance comparison and direct concentration}\label{sec:variance}

Let $\mathcal C$ be the countable collection of positive-probability batch
types. Independent Poisson processes $P_C$ of rates $\mu(C)$ realize the
Poissonized experiment. Define
\[
 E_C=\inf\{s:P_C(s)>0\},\qquad T_i=\inf_{C\ni i}E_C.
\]
The $E_C$ are independent, but the $T_i$ generally are not. For each $i$,
$T_i\sim\Exp(\lambda_i)$, and
\begin{equation}\label{eq:jointabsence}
 \Pp(T_i>t,T_j>t)
 =e^{-t(\lambda_i+\lambda_j-\lambda_{ij})}.
\end{equation}
Consequently,
\begin{equation}\label{eq:exactvariance}
 V(t)=\sum_{i,j}\left[
 e^{-t(\lambda_i+\lambda_j-\lambda_{ij})}
 -e^{-t(\lambda_i+\lambda_j)}\right].
\end{equation}
The diagonal convention in \eqref{eq:rates} makes this formula valid also for
$i=j$.

\begin{theorem}[Sharp variance comparison]\label{thm:varcomparison}
For every $t\ge0$,
\begin{equation}\label{eq:sharpcomparison}
 v_0(t)\le V(t)\le Bv_0(t).
\end{equation}
Both constants are optimal over the class of laws with batch size at most $B$.
For arbitrary real weights $(w_i)$, any well-defined $L^2$ weighted sum satisfies
\begin{equation}\label{eq:signedweightedvariance}
 \Var\left(\sum_iw_i\ind_{\{T_i\le t\}}\right)
 \le B\sum_iw_i^2e^{-t\lambda_i}(1-e^{-t\lambda_i}).
\end{equation}
\end{theorem}
\begin{proof}
Every term of \eqref{eq:exactvariance} is nonnegative; its diagonal sum is
$v_0(t)$. For the upper bound, each independent clock $E_C$ is read by at most
$|C|\le B$ of the indicators. Apply \eqref{eq:readvar}, also with the real
weights absorbed into the individual functions. Countable sums converge in
$L^2$ because $X(t)\le BN(t)$, or by the centered-series version of that bound.
Singleton batches attain the lower constant. If batch types are disjoint
blocks of exactly $B$ species, their discovery indicators are independent
between blocks and identical within each block; then $V(t)=Bv_0(t)$.
\end{proof}

\begin{proposition}[Species-dependent variance bound]\label{prop:speciesvar}
A deterministic cardinality bound is unnecessary for the following refinement.
For a law on finite batches, set
\[
 \kappa_i=\E[|A_1|\mid i\in A_1],\qquad
 v_i(t)=e^{-t\lambda_i}(1-e^{-t\lambda_i}).
\]
For every finite set $J$ of species and real weights $w_i$,
\begin{equation}\label{eq:speciesvar}
 \Var\left(\sum_{i\in J}w_i\ind_{\{T_i\le t\}}\right)
 \le\sum_{i\in J}\kappa_iw_i^2v_i(t).
\end{equation}
Infinite values on the right give a vacuous inequality. If the full right-hand
sum is finite, the corresponding centered series converges in $L^2$ and
\eqref{eq:speciesvar} holds for that limit. An uncentered series must additionally
be well-defined. Under \eqref{eq:batchbound}, $\kappa_i\le B$.
\end{proposition}
\begin{proof}
The finite weighted sum has a terminal Doob martingale whose jump at a batch
$C$ at time $s$ is $\sum_{i\in C\cap J}w_i g_i(s)$, where
$g_i(s)=e^{-\lambda_i(t-s)}\ind_{\{T_i>s\}}$ (with left limits at arrivals).
This conditional-expectation calculation is the weighted version of
\eqref{eq:H}. The jumps are bounded by $\sum_{i\in J}|w_i|$, irrespective of
$|C|$, so martingale isometry applies. By Cauchy--Schwarz within each batch,
\begin{align*}
 \Var\left(\sum_{i\in J}w_i\ind_{\{T_i\le t\}}\right)
 &\le \E\int_0^t\sum_C\mu(C)|C|
            \sum_{i\in C\cap J}w_i^2g_i(s)^2\,ds\\
 &=\sum_{i\in J}\kappa_i w_i^2\int_0^t
       \lambda_i e^{-2\lambda_i(t-s)}e^{-\lambda_i s}\,ds
 =\sum_{i\in J}\kappa_iw_i^2v_i(t).
\end{align*}
Tonelli justifies the nonnegative sums. Apply the finite inequality to tails
to obtain the stated centered $L^2$ extension. Equivalently, the key identity
is $\sum_j\lambda_{ij}=\kappa_i\lambda_i$.
\end{proof}

\begin{theorem}[Read-$B$ Bennett concentration]\label{thm:bennett}
For $t>0$ with $v_0(t)>0$, and $x>0$,
\begin{align}
 \Pp(|X(t)-m(t)|\ge x)
 &\le2\exp\left\{-\frac{v_0(t)}B h\left(\frac{x}{v_0(t)}\right)\right\}
 \label{eq:discoverybennett}\\
 &\le2\exp\left\{-\frac{x^2}{2Bv_0(t)+(2B/3)x}\right\}.
 \label{eq:discoverybernstein}
\end{align}
\end{theorem}
\begin{proof}
Apply Corollary~\ref{cor:readB} and use
$h(u)\ge u^2/(2(1+u/3))$ for $u\ge0$.
\end{proof}

The factor $Bv_0(t)$ in the quadratic part of this moment bound cannot be
uniformly reduced: its second derivative at zero must dominate the true
variance, and the disjoint-block example attains $V(t)=Bv_0(t)$. This is a
sharpness statement about the universal variance proxy, not a claim that every
numerical constant in a tail inequality is optimal.

\begin{proposition}[Fractional, species-dependent Bennett bound]\label{prop:fractionalbennett}
Suppose $\E|A_1|<\infty$ and choose $\alpha_i>0$ such that
\begin{equation}\label{eq:fractionalcover}
 \sum_{i\in C}\alpha_i\le1\quad\text{whenever }\mu(C)>0.
\end{equation}
Then, with the right side allowed to be infinite,
\begin{equation}\label{eq:fractionalmgf}
 \log\E e^{\theta(X(t)-m(t))}
 \le\sum_i\alpha_i v_i(t)\psi(|\theta|/\alpha_i).
\end{equation}
In particular, if $B_i=\sup\{|C|:\mu(C)>0,\ i\in C\}<\infty$ for each $i$,
then $\alpha_i=1/B_i$ is admissible and
\begin{equation}\label{eq:speciesbennett}
 \log\E e^{\theta(X(t)-m(t))}
 \le\sum_i\frac{v_i(t)}{B_i}\psi(B_i|\theta|).
\end{equation}
\end{proposition}
\begin{proof}
Apply Lemma~\ref{lem:finner} to a finite set of centered discovery indicators,
with weights $\alpha_i$ and functions
$\exp\{\theta(\ind_{\{T_i\le t\}}-\Pp(T_i\le t))/\alpha_i\}$.
The Bernoulli moment estimate in Corollary~\ref{cor:readB} gives
\eqref{eq:fractionalmgf}. Since $\E X(t)\le t\E|A_1|<\infty$, the centered
partial sums converge almost surely to $X(t)-m(t)$; Fatou's lemma proves
the countable assertion. For the special choice, $B_i\ge|C|$ for $i\in C$,
so $\sum_{i\in C}1/B_i\le1$.
\end{proof}

The function $b\mapsto\psi(bu)/b$ is increasing for $b>0$, $u\ge0$.
Thus \eqref{eq:speciesbennett} improves the global-$B$ moment bound whenever
some $B_i<B$. This is a \emph{support-sensitive}, not a probability-sensitive,
improvement: a very rare large batch still enlarges every $B_i$ on that batch.
In contrast, $\kappa_i$ in Proposition~\ref{prop:speciesvar} averages batch size
conditional on containing $i$. More general feasible choices in
\eqref{eq:fractionalcover} may improve on $1/B_i$. Chernoff optimization of
\eqref{eq:fractionalmgf} supplies the associated two-sided tail bound.

More generally, the independent coordinates can be the entire paths
$P_C|_{[0,t]}$. For species-wise functions of their incident paths,
\eqref{eq:readvar} and \eqref{eq:readmgf} still hold. In particular,
\begin{equation}\label{eq:generalfunctionalvariance}
 \Var\left(\sum_i g_i(N_i(t))\right)
 \le B\sum_i\Var(g_i(\Poi(t\lambda_i))),\qquad
 N_i(t)=\sum_{C\ni i}P_C(t),
\end{equation}
under the stated convergence requirements. No monotonicity is needed for this
upper bound. A covariance lower bound, in contrast, requires additional sign
or monotonicity information.

\section{The terminal martingale and its quadratic variation}\label{sec:bracket}

Fix a terminal time $t$. Let $\mathcal F_s$ record all batch arrivals up to $s$,
and put
\[
 M_s=\E[X(t)\mid\mathcal F_s]-m(t),\qquad 0\le s\le t.
\]
With $Y_i(s)=\ind_{\{T_i>s\}}$, an explicit conditional expectation is
\[
 \E[X(t)\mid\mathcal F_s]
 =X(s)+\sum_iY_i(s)(1-e^{-\lambda_i(t-s)}).
\]
The sum is finite in expectation by \eqref{eq:rows}. At an arrival of batch
$C$, the martingale jump is
\begin{equation}\label{eq:H}
 H_C(s)=\sum_{i\in C}e^{-\lambda_i(t-s)}Y_i(s-),
 \qquad 0\le H_C(s)\le B.
\end{equation}
Thus
\begin{equation}\label{eq:martrep}
 M_s=\sum_C\int_0^s H_C(u)\{dP_C(u)-\mu(C)\,du\},
\end{equation}
and
\begin{equation}\label{eq:bracketdef}
 Q_t:=\langle M\rangle_t
 =\int_0^t\sum_C\mu(C)H_C(s)^2\,ds,
 \qquad \E Q_t=V(t).
\end{equation}
All integrands are bounded, the total event rate is one, and $Q_t\le B^2t$.
In particular, the stochastic integrals are square-integrable.

\begin{theorem}[Positive decomposition and bounded bracket overlap]\label{thm:bracket}
Let
\begin{equation}\label{eq:LB}
 L_B=\frac{B(4B-3)}2\le2B^2.
\end{equation}
Then
\begin{equation}\label{eq:bracketvar}
 \Var(Q_t)\le L_BV(t),
\end{equation}
and, for every $\theta\in\R$,
\begin{equation}\label{eq:bracketmgf}
 \log\E e^{\theta(Q_t-V(t))}
 \le\frac{V(t)}{L_B}\psi(L_B\theta).
\end{equation}
In particular $L_2=5$ is valid for dependent pairs.
\end{theorem}
\begin{proof}
Expanding \eqref{eq:bracketdef} gives
\begin{equation}\label{eq:bracketexpanded}
 Q_t=\sum_{i,j}\lambda_{ij}\int_0^t
 e^{-(\lambda_i+\lambda_j)(t-s)}
 \ind_{\{T_i>s,T_j>s\}}\,ds.
\end{equation}
Set $r_{ij}=\lambda_i+\lambda_j$, $a_{ij}=\lambda_{ij}/r_{ij}$, and
\begin{equation}\label{eq:Fij}
 F_{ij}=e^{-r_{ij}(t-\min\{t,T_i,T_j\})}-e^{-r_{ij}t}.
\end{equation}
Direct integration proves
\begin{equation}\label{eq:Qdecomp}
 Q_t=\sum_{i,j}a_{ij}F_{ij},\qquad 0\le F_{ij}\le1.
\end{equation}
The support of $F_{ij}$ consists of the batch clocks meeting $\{i,j\}$.
For a batch $C$, its total coefficient load is
\begin{align*}
 \sum_{i\in C\ \text{or}\ j\in C}a_{ij}
 &=2\sum_{i\in C}\sum_j a_{ij}-\sum_{i,j\in C}a_{ij}.
\end{align*}
Since $a_{ii}=1/2$ and
\[
 \sum_{j\ne i}a_{ij}
 \le\frac1{\lambda_i}\sum_{j\ne i}\lambda_{ij}\le B-1,
\]
this load is at most
\[
 2|C|(B-1/2)-|C|/2=|C|(2B-3/2)\le L_B.
\]
Apply Lemma~\ref{lem:variance-drop} and
Proposition~\ref{prop:overlapmgf}, using $\E Q_t=V(t)$.
The argument is valid for countably many pairs by their nonnegative partial
sums. Each partial sum is bounded by $Q_t\le B^2t$, so all required limits
also follow directly by dominated convergence.
\end{proof}

\begin{proposition}[Sharpness of the coefficient-load constant]\label{prop:loadsharp}
For the coefficients $a_{ij}=\lambda_{ij}/(\lambda_i+\lambda_j)$ and the
supports in \eqref{eq:Qdecomp}, the universal upper bound $L_B$ on the load
of one batch clock cannot be reduced.
\end{proposition}
\begin{proof}
For $B=1$ the load is exactly $a_{ii}=1/2=L_1$. For $B\ge2$, take hubs
$1,\ldots,B$ and a central batch $C_0=\{1,\ldots,B\}$. For each hub add
$N$ batches containing that hub and $B-1$ fresh leaves, all leaves distinct.
Give the $BN+1$ batches equal probability. Each hub has marginal rate
$(N+1)/(BN+1)$; each leaf has rate $1/(BN+1)$. The load of $C_0$ is
\[
 \frac B2+\frac{2BN(B-1)}{N+2}
             +\frac{B(B-1)}{2(N+1)}
 \longrightarrow \frac B2+2B(B-1)=L_B.
\]
The three terms come from the hub diagonals, hub--leaf ordered pairs, and
distinct ordered hub pairs, respectively.
\end{proof}

This sharpness concerns the \emph{worst-case coefficient load}, not the best
possible constant in $\Var(Q_t)\le c_B\E Q_t$. Exploiting the particular
functions $F_{ij}$ or changing the representation could still improve $c_B$.

\begin{corollary}[Bracket tails and moments]\label{cor:brackettails}
For $\delta>0$,
\begin{equation}\label{eq:Qupper}
 \Pp(Q_t\ge(1+\delta)V(t))
 \le\exp\{-V(t)h(\delta)/L_B\}.
\end{equation}
For $0<\delta<1$,
\begin{equation}\label{eq:Qlower}
 \Pp(Q_t\le(1-\delta)V(t))
 \le\exp\{-V(t)h(-\delta)/L_B\}.
\end{equation}
For every fixed $p\ge2$,
\begin{equation}\label{eq:Qmoments}
 \E|Q_t-V(t)|^p
 \le C_p\{(B^2V(t))^{p/2}+B^{2p}\}.
\end{equation}
\end{corollary}
\begin{proof}
Chernoff optimization gives the two tails. The moment estimate follows by
integrating the resulting Bernstein bound, with variance proxy $L_BV(t)$ and
scale $L_B$, and using $L_B\le2B^2$.
\end{proof}

\section{Quantitative normal approximation and triangular arrays}\label{sec:normal}

For a real random variable $Z$, write
$\dK(Z,\Normal)=\sup_x|\Pp(Z\le x)-\Phi(x)|$.
We use the following standard martingale normal-approximation inequality:
for martingale differences $(\zeta_k)$ with
$\E(\sum_k\zeta_k)^2=1$, and $p\ge1$,
\begin{equation}\label{eq:Haeusler}
 \dK\left(\sum_k\zeta_k,\Normal\right)
 \le C_p\left\{
 \E\left|\sum_k\E[\zeta_k^2\mid\mathcal G_{k-1}]-1\right|^p
 +\sum_k\E|\zeta_k|^{2p}\right\}^{1/(2p+1)}.
\end{equation}
This is Haeusler's inequality~\cite{Haeusler}, in the formulation recorded in
Mourrat~\cite[Theorem~1.1]{Mourrat}. That inequality requires the displayed
moments, not bounded increments. We use fixed orders $p\ge2$ below.
Appendix~\ref{app:mesh} proves the exact
deterministic-grid passage needed for our finite-activity continuous-time
martingales; no independent-increment assertion about $M$ is used.

\begin{theorem}[Uniform Poissonized Berry--Esseen bounds]\label{thm:BE}
There is a universal $C<\infty$ such that, for every bounded-batch law and every
$t$ with $V(t)>0$,
\begin{equation}\label{eq:BE}
 \dK\left(\frac{X(t)-m(t)}{\sqrt{V(t)}},\Normal\right)
 \le \min\left\{1,C\left(\frac{B^2}{V(t)}\right)^{1/5}\right\}.
\end{equation}
For each fixed $p\ge2$, if $V(t)\ge B^2$, then
\begin{equation}\label{eq:BEp}
 \dK\left(\frac{X(t)-m(t)}{\sqrt{V(t)}},\Normal\right)
 \le C_p\left(\frac{B^2}{V(t)}\right)^{p/(4p+2)}.
\end{equation}
\end{theorem}
\begin{proof}
Let $V=V(t)$. The normalized terminal martingale has bracket $Q_t/V$. Moreover,
for $p\ge2$, the compensator identity and \eqref{eq:H} give
\begin{align}
 \E\sum_{s\le t}|\Delta M_s|^{2p}
 &=\E\int_0^t\sum_C\mu(C)H_C(s)^{2p}\,ds
 \le B^{2p-2}V.\label{eq:jumpmoments}
\end{align}
Apply the continuous-time consequence of \eqref{eq:Haeusler} from
Appendix~\ref{app:mesh}. For $p=2$, \eqref{eq:bracketvar} and
\eqref{eq:jumpmoments} give a bound by $C(B^2/V)^{1/5}$.
For general $p$, put $\eta=B^2/V\le1$. Equations~\eqref{eq:Qmoments} and
\eqref{eq:jumpmoments} bound the quantity in braces by
$C_p(\eta^{p/2}+\eta^p+\eta^{p-1})\le C'_p\eta^{p/2}$.
This proves \eqref{eq:BEp}.
\end{proof}

\begin{corollary}[Changing populations and growing batch bounds]\label{cor:triangular}
Consider any sequence of Poissonized models, with possibly changing laws,
observation times, and batch bounds $B_k$. Write $X_k$ for the discovery count
and $V_k=\Var(X_k)>0$. If
\begin{equation}\label{eq:triangular}
 B_k/\sqrt{V_k}\longrightarrow0,
\end{equation}
then $(X_k-\E X_k)/\sqrt{V_k}\dto\Normal$.
\end{corollary}
\begin{proof}
Apply \eqref{eq:BE} model by model. In particular, for a fixed $B$ this proves
\eqref{eq:poisCLT}.
\end{proof}

For fixed $B$, \eqref{eq:BEp} implies
$\dK=O_\varepsilon(V^{-1/4+\varepsilon})$ for every $\varepsilon>0$.
No optimality claim is made for this rate in the discovery class. The constants
$C_p$ depend on $p$; one may not send $p$ to infinity without controlling them.

\subsection{Stopping, Gaussian completion, and a logarithmic refinement}

The next argument improves the fixed-moment rates without increasing $p$.
Stopping and augmentation to a deterministic bracket are standard
martingale-approximation devices (cf.\ \cite[proof of Theorem~1.5]{Mourrat}); we give the completion and smoothing
steps needed here in full. Only \eqref{eq:Haeusler} at $p=3$ is used as an
external normal-approximation theorem.

\begin{lemma}[Removal of a conditionally Gaussian perturbation]\label{lem:smoothing}
Let $Z$ and $\sigma\ge0$ be jointly distributed, and let $G\sim\Normal$ be
independent of $(Z,\sigma)$. Put $Y=Z+\sigma G$. If $a>0$ and
$\Pp(\sigma>a)\le\rho$, then
\begin{equation}\label{eq:smoothing}
 \dK(Z,\Normal)\le C\{\dK(Y,\Normal)+a+\rho\}
\end{equation}
for a universal constant $C$.
\end{lemma}
\begin{proof}
Put $d=\dK(Y,\Normal)$ and $A=\{\sigma\le a\}$. For any interval $I$ of
length $a$, independence of $G$ implies
\[
 \Pp(Z\in I,A)\Pp(|G|\le1)
 \le\Pp(Y\in I+[-a,a])\le3a/\sqrt{2\pi}+2d.
\]
The bound also holds for closed or half-open intervals by taking one-sided
limits of the distribution functions. Fix $x$ and partition each side of
$x$ into intervals of length $a$, starting at distance $ka$, $k=0,1,\ldots$.
On $A$, conditional on $(Z,\sigma)$ in such an interval, the probability
that $Z$ and $Y$ lie on different sides of $x$ is at most $1-\Phi(k)$.
The sum of these Gaussian tails is finite. Hence
\[
 |\Pp(Z\le x)-\Pp(Y\le x)|\le C(a+d)+\rho.
\]
Take the supremum and use the triangle inequality with the normal law.
No independence between $Z$ and $\sigma$ is required.
\end{proof}

\begin{lemma}[Completion to deterministic bracket]\label{lem:completion}
Let $M$ be the discovery martingale stopped at
\[
 \tau=\inf\{s\le t:Q_s\ge\nu\}\wedge t,
 \qquad \nu>0,
\]
where an empty infimum is interpreted as $+\infty$. Set
$Q_\tau=\min(Q_t,\nu)$ and let $G\sim\Normal$ be independent of
$\mathcal F_t$. Then
\begin{equation}\label{eq:completion}
 \dK\left(\frac{M_\tau+\sqrt{\nu-Q_\tau}\,G}{\sqrt\nu},\Normal\right)
 \le C\left(\frac{B^4}{\nu^2}\right)^{1/7}.
\end{equation}
The same statement holds for the threshold and weighted martingales, with
their respective jump bounds.
\end{lemma}
\begin{proof}
Continuity of the predictable bracket gives $Q_\tau=\min(Q_t,\nu)$,
with no overshoot. Discretize $M_{s\wedge\tau}$ on deterministic time grids.
Its grid predictable variance $Q^\pi$ tends to $Q_\tau$ in $L^3$ and its
summed sixth moments tend to $\E\sum_{s\le\tau}|\Delta M_s|^6$; the stopped
version of the grid argument is proved in Appendix~\ref{app:mesh}.
After the grid increments, append $m$ independent standard Gaussian variables
$G_1,\ldots,G_m$, independent of $\mathcal F_t$, multiplied by
$\sqrt{(\nu-Q_\tau)/m}$. Enlarge the filtration at the end of the grid to
$\mathcal F_t$ (a zero martingale difference may be inserted there).
The appended conditional variances sum to $\nu-Q_\tau$.
The terminal distribution is that on the left of \eqref{eq:completion}, for
every $m$ and every grid, and its unconditional variance is exactly $\nu$.
The normalized predictable variance differs from one by
$(Q^\pi-Q_\tau)/\nu$. Moreover,
\[
 \E\sum_{s\le\tau}|\Delta M_s|^6\le B^4\E Q_\tau\le B^4\nu,
 \qquad
 \sum_{j=1}^m\E\left|\sqrt{\frac{\nu-Q_\tau}{m}}G_j\right|^6
 =\frac{15}{m^2}\E(\nu-Q_\tau)^3.
\]
Apply \eqref{eq:Haeusler} at $p=3$, first let the grid mesh tend to zero,
and then let $m\to\infty$. This proves \eqref{eq:completion}.
\end{proof}

\begin{theorem}[Logarithmically improved normal approximation]\label{thm:BElog}
There is a universal constant $C$ such that, with $\eta=B^2/V(t)>0$,
\begin{equation}\label{eq:BElog}
 \dK\left(\frac{X(t)-m(t)}{\sqrt{V(t)}},\Normal\right)
 \le \mathcal R(\eta),\qquad
 \mathcal R(\eta)=\min\{1,C[\eta\log(e+\eta^{-1})]^{1/4}\}.
\end{equation}
The constant can be chosen once for all laws, times, and batch bounds.
\end{theorem}
\begin{proof}
Write $V=V(t)$ and take $0<\delta<1/2$. Set $\nu=(1+\delta)V$ and use
Lemma~\ref{lem:completion}. Changing the normalizing scale from $\sqrt\nu$
to $\sqrt V$ costs at most $C\delta$, since
$\sup_x|\Phi(x/\sqrt{1+\delta})-\Phi(x)|\le C\delta$.
For $Z=M_\tau/\sqrt V$ and $\sigma=\sqrt{(\nu-Q_\tau)/V}$,
$\sigma\le\sqrt{2\delta}$ outside $\{Q_t<(1-\delta)V\}$.
Also $M_\tau=M_t$ outside $\{Q_t\ge(1+\delta)V\}$.
The bracket tails and $L_B\le2B^2$ therefore give, with universal $c>0$,
\begin{equation}\label{eq:BElog-intermediate}
 \dK(M_t/\sqrt V,\Normal)
 \le C\left\{\eta^{2/7}+\sqrt\delta+
                      e^{-c\delta^2/\eta}\right\}.
\end{equation}
Here Lemma~\ref{lem:smoothing} removes the added Gaussian. For sufficiently
small $\eta$, choose $\delta=A\sqrt{\eta\log(1/\eta)}$, with a fixed
sufficiently large $A$. Then the right side of
\eqref{eq:BElog-intermediate} is at most
$C[\eta\log(1/\eta)]^{1/4}$. In particular,
$\eta^{2/7}=o([\eta\log(1/\eta)]^{1/4})$.
Increase $C$ and use $\dK\le1$ for the remaining values of $\eta$.
\end{proof}

\section{An absolute de-Poissonization variance expansion}\label{sec:correction}

The next statement does not need a deterministic bound on batch size.
It is a fixed-law assertion: its $o(1)$ is not claimed uniformly over changing
sampling distributions.

\begin{theorem}[Universal fixed-law variance correction]\label{thm:correction}
Suppose the batch law is fixed and $\E|A_1|^2<\infty$. With the notation
\eqref{eq:means}--\eqref{eq:flux},
\begin{equation}\label{eq:meancorrection}
 b_n-m(n)\longrightarrow0,
\end{equation}
and
\begin{equation}\label{eq:variancecorrection}
 s_n^2=V(n)-nf(n)^2+o(1).
\end{equation}
Thus the first correction caused by fixing the number of batches depends
only on the marginal rates $\lambda_i$.
\end{theorem}

\begin{lemma}[A uniformly Lipschitz binomial--exponential discrepancy]\label{lem:delta}
For $n\ge1$ and $0\le x\le1$, let
\[
 \delta_n(x)=e^{-nx}-(1-x)^n.
\]
Then $0\le\delta_n(x)\le2x$, and
\begin{equation}\label{eq:deltabound}
 \sup_{n\ge1}\sup_{0\le x\le1}|\delta_n'(x)|\le2.
\end{equation}
For every $\epsilon>0$,
$\sup_{\epsilon\le x\le1}|\delta_n'(x)|\to0$.
\end{lemma}
\begin{proof}
For $x\le1/2$, use
$0\le-\log(1-x)-x\le x^2$ to obtain
$\delta_n(x)\le nx^2e^{-nx}\le x/e$. For $x>1/2$ the bound
$\delta_n(x)\le1\le2x$ suffices.
For $n\ge2$,
\[
 \delta_n'(x)=n\{(1-x)^{n-1}-e^{-nx}\}.
\]
On $[0,1/2]$, insert $e^{-(n-1)x}$ between the two terms to obtain
\[
 |\delta_n'(x)|
 \le n(n-1)x^2e^{-(n-1)x}+nxe^{-(n-1)x}
 \le2(4e^{-2}+e^{-1})<2.
\]
For $x>1/2$, the bound is
$n2^{-(n-1)}+ne^{-n/2}\le1+2/e<2$.
The case $n=1$ is immediate. Finally, on $[\epsilon,1]$,
$|\delta_n'(x)|\le n\{(1-\epsilon)^{n-1}+e^{-n\epsilon}\}\to0$.
\end{proof}

\begin{proof}[Proof of Theorem~\ref{thm:correction}]
The means differ by $\sum_i\delta_n(\lambda_i)$. The summands converge to zero
and are bounded by $2\lambda_i$, whose sum is finite. This proves
\eqref{eq:meancorrection}.

For each ordered pair, abbreviate
\[
 a=\lambda_i,\quad b=\lambda_j,\quad c=\lambda_{ij},\qquad
 u=a+b-c,\quad w=a+b-ab.
\]
Both $u$ and $w$ lie in $[0,1]$ and are at least $\max(a,b)$.
The exact fixed-sample covariance, including when $i=j$, is
\[
 (1-u)^n-(1-a)^n(1-b)^n=(1-u)^n-(1-w)^n.
\]
It follows that
\begin{equation}\label{eq:covsplit}
 s_n^2=V(n)+E_n-R_n,
\end{equation}
where
\begin{align}
 E_n&=\sum_{i,j}\{\delta_n(w)-\delta_n(u)\},\label{eq:En}\\
 R_n&=\sum_{i,j}e^{-n(a+b)}(e^{nab}-1).\label{eq:Rn}
\end{align}
These manipulations are justified by absolute summability: the counts have
finite second moments, and the estimates below also dominate the discrepancy
terms explicitly.

By Lemma~\ref{lem:delta},
\[
 |\delta_n(w)-\delta_n(u)|\le2|c-ab|\le2(c+ab).
\]
For each fixed pair of positive-rate species the left side tends to zero,
because the interval between $u$ and $w$ stays away from zero. Moreover,
\[
 \sum_{i,j}c=\E|A_1|^2<\infty,
 \qquad \sum_{i,j}ab=(\E|A_1|)^2<\infty.
\]
Thus $E_n\to0$ by dominated convergence. Notice that this argument never
introduces $(1-a-b)^n$ when $a+b>1$.

The linear term in \eqref{eq:Rn} is exactly $nf(n)^2$.
Since $e^z-1-z\le z^2e^z/2$ for $z\ge0$, and
$w=a+b-ab\ge(a+b)/2$,
\begin{align}
 0\le R_n-nf(n)^2
 &\le\frac{n^2}{2}\sum_{i,j}a^2b^2e^{-nw}\notag\\
 &\le\frac12\left(n\sum_i\lambda_i^2e^{-n\lambda_i/2}\right)^2=o(1).
 \label{eq:Rremainder}
\end{align}
For the last step, write the inner sum as
$\sum_i\lambda_i(n\lambda_i)e^{-n\lambda_i/2}$ and dominate it by
$(2/e)\sum_i\lambda_i$. Combining this with \eqref{eq:covsplit} proves
\eqref{eq:variancecorrection}. No diagonal terms have been discarded.
\end{proof}

\begin{remark}
The first-order term $nf(n)^2$ need not be bounded. Equation
\eqref{eq:variancecorrection} is an \emph{absolute} $o(1)$ remainder after
subtracting it, not a claim that the two variances differ by $o(1)$.
The fixed-law restriction is material. For example, for $n$ equally probable
species and $n$ draws, the difference between fixed and Poissonized means
tends to $1/(2e)$, rather than zero, as the law changes with $n$.
\end{remark}

\section{Fixed-sample CLT and quantitative transfer}\label{sec:depois}

\begin{lemma}[Discovery-derivative estimate]\label{lem:flux}
For a fixed summable sequence of positive rates with infinitely many terms,
\begin{equation}\label{eq:fluxratio}
 \frac{f(t)^2}{f(2t)}\longrightarrow0,
 \qquad tf(t)^2=o(v_0(t)).
\end{equation}
\end{lemma}
\begin{proof}
For any finite set $F$ of species, its contribution to $f(t)$ is negligible
relative to $f(t)$: for each $i\in F$ choose a positive rate $\lambda_j<\lambda_i$
and compare $\lambda_i e^{-t\lambda_i}$ with
$\lambda_j e^{-t\lambda_j}$. Consequently,
\[
 f(t)=(1+o(1))\sum_{i\notin F}\lambda_i e^{-t\lambda_i}.
\]
Cauchy--Schwarz yields
\[
 \limsup_{t\to\infty}\frac{f(t)^2}{f(2t)}
 \le\sum_{i\notin F}\lambda_i.
\]
Let $F$ increase to all species. Finally,
\[
 v_0(t)=m(2t)-m(t)=\int_t^{2t}f(u)\,du\ge tf(2t).
\]
This is the infinite-urn estimate used in
\cite{Dutko,ErikssonGraphs}, written in discovery-derivative notation.
\end{proof}

\begin{proposition}[Variance equivalence]\label{prop:variance-equivalence}
For a fixed law with $\E|A_1|^2<\infty$, if either $s_n^2$ or $V(n)$ tends to infinity,
then both do, and $s_n^2\sim V(n)$.
\end{proposition}
\begin{proof}
Divergence of either variance forces infinitely many positive-rate species.
The exact covariance formula and its nonnegative off-diagonal terms still
give $v_0(n)\le V(n)$ under this moment assumption. By Lemma~\ref{lem:flux},
$nf(n)^2=o(v_0(n))=o(V(n))$. Theorem~\ref{thm:correction} now gives
$s_n^2=V(n)(1+o(1))+o(1)$. Also $s_n^2\le V(n)+o(1)$ directly from
\eqref{eq:variancecorrection}, so $s_n^2\to\infty$ forces $V(n)\to\infty$.
\end{proof}

\begin{lemma}[A square-root window is negligible]\label{lem:window}
Suppose the law is fixed, $\E|A_1|<\infty$, each $V(n)$ is finite, and
$V(n)\to\infty$. For fixed $C>0$ and
$h_n=\lceil C\sqrt n\rceil$,
\begin{equation}\label{eq:window}
 b_{n+h_n}-b_{n-h_n}=o(\sqrt{V(n)}).
\end{equation}
\end{lemma}
\begin{proof}
For all sufficiently large $n$, $n-h_n\ge1$. The mean comparison gives
\[
 b_{n+h_n}-b_{n-h_n}
 =m(n+h_n)-m(n-h_n)+o(1)\le2h_nf(n-h_n)+o(1).
\]
Set $\delta_n=n^{-2/3}$. Split the rates at $\delta_n$. Then
\[
 f(n-h_n)\le e^{h_n\delta_n}f(n)
       +\E|A_1|\, e^{-(n-h_n)\delta_n}.
\]
Here $h_n\delta_n\to0$, the second term decays faster than any negative
power of $n$, and $\sqrt n f(n)=o(\sqrt{v_0(n)})$ by
Lemma~\ref{lem:flux}. Since $v_0(n)\le V(n)$ and $V(n)\to\infty$,
\eqref{eq:window} follows.
\end{proof}

\begin{proof}[Completion of Theorem~\ref{thm:main}]
Couple $D_n$ and $X(n)=D_{N(n)}$ using the same batch sequence. On
$\{|N(n)-n|\le h_n\}$,
\[
 |D_{N(n)}-D_n|\le D_{n+h_n}-D_{n-h_n}.
\]
For every $\epsilon>0$, Markov's inequality and
Lemma~\ref{lem:window} show that
\[
 \limsup_{n\to\infty}
 \Pp\bigl(|D_{N(n)}-D_n|>\epsilon\sqrt{V(n)}\bigr)\le C^{-2},
\]
using $\Var(N(n))=n$. Send $C\to\infty$. Thus
\begin{equation}\label{eq:coupling-negligible}
 D_{N(n)}-D_n=o_{\Pp}(\sqrt{V(n)}).
\end{equation}
The Poissonized CLT, \eqref{eq:meancorrection}, and
Proposition~\ref{prop:variance-equivalence} imply \eqref{eq:fixedCLT} by
Slutsky's theorem. The equivalence with $v_0(n)\to\infty$ follows from
\eqref{eq:sharpcomparison}.
\end{proof}

\subsection{An explicit transfer bound}
The qualitative transfer above does not assert a fixed-sample error bound
in terms of variance alone. The following finite-sample inequality makes the
additional terms visible.

\begin{proposition}[Quantitative de-Poissonization]\label{prop:quantdepois}
Let $V=V(n)>0$, $1\le h\le n$, and $\epsilon>0$. Then
\begin{align}
 \dK\left(\frac{D_n-b_n}{\sqrt V},\Normal\right)
 \le{}&\mathcal R(B^2/V)
 +\frac{\epsilon+|b_n-m(n)|/\sqrt V}{\sqrt{2\pi}}\notag\\
 &+2\exp\left\{-\frac{h^2}{2(n+h/3)}\right\}
 +\frac{b_{n+h}-b_{n-h}}{\epsilon\sqrt V}.
 \label{eq:quantdepois}
\end{align}
\end{proposition}
\begin{proof}
Under the preceding coupling, except on an event of probability at most
\[
 \Pp(|N(n)-n|>h)
 +\frac{b_{n+h}-b_{n-h}}{\epsilon\sqrt V},
\]
the two centered normalized counts differ by at most
$\epsilon+|b_n-m(n)|/\sqrt V$. Apply the bound on the normal density,
Theorem~\ref{thm:BElog}, and the Poisson Bernstein inequality
$\Pp(|N(n)-n|>h)\le2e^{-h^2/(2(n+h/3))}$.
\end{proof}

Replacing $\sqrt V$ by $s_n$ adds a normal scale-comparison error of order
$|s_n^2/V-1|$ when that ratio is bounded away from zero and infinity.
The expansion in Section~\ref{sec:correction} controls this error for a fixed
law. It does not eliminate the discovery-window term uniformly over all laws.

\subsection{Fixed-sample mean-based concentration}

\begin{proposition}[Self-bounding concentration]\label{prop:selfbound}
For any bounded-batch law, any $n$, and $b_n>0$,
\begin{align}
 \Pp(D_n-b_n\ge x)
 &\le\exp\left\{-\frac{b_n}B h\left(\frac{x}{b_n}\right)\right\},\label{eq:fixed-upper}\\
 \Pp(D_n-b_n\le-x)
 &\le\exp\left\{-\frac{x^2}{2Bb_n}\right\},\qquad x>0.\label{eq:fixed-lower}
\end{align}
\end{proposition}
\begin{proof}
Remove batch $k$ and call the resulting count $D_n^{(-k)}$. Then
\[
 0\le D_n-D_n^{(-k)}\le B,\qquad
 \sum_{k=1}^n(D_n-D_n^{(-k)})\le D_n.
\]
A species contributes to this sum only if it occurs in exactly one of the
batches. Hence $D_n/B$ is self-bounding, and the standard self-bounding
exponential-moment inequalities~\cite[p.~1885]{BLM} give the Bennett upper bound. More specifically, the
lower bound is also the final lower-tail inequality in
\cite[Theorem~1]{BLM}, with $(a,b)=(1,0)$, applied to $D_n/B$.
It is stated there for $0<x<b_n$; the endpoint follows by continuity from
below, and $x>b_n$ is impossible. The full sub-Poisson moment inequality is
originally due to Boucheron, Lugosi, and Massart~\cite{BLM2000}.
\end{proof}

The read-$B$ argument cannot simply be applied to the original sampling
coordinates $A_k$: as a function of a possible batch value, one such coordinate
can affect many different species indicators. The read-$B$ property applies to
the \emph{independent Poisson batch-type coordinates}. This is why the fixed-
sample bounds here are stated separately and use $Bb_n$, not the actual variance.

\section{Unbounded batches: a Lindeberg extension and an obstruction}\label{sec:unbounded}

\begin{theorem}[Poissonized Lindeberg extension]\label{thm:unbounded}
Consider a sequence of laws on finite subsets, observation times $t_\ell$,
and batch sizes $J_\ell=|A_{\ell,1}|$. Assume $\E J_\ell^2<\infty$ for each
$\ell$. Let $X_\ell$ be the Poissonized discovery count and suppose
$V_\ell=\Var(X_\ell)\to\infty$. If, for every $\epsilon>0$,
\begin{equation}\label{eq:batchlindeberg}
 t_\ell\E[J_\ell^2;J_\ell>\epsilon\sqrt{V_\ell}]=o(V_\ell),
\end{equation}
then $(X_\ell-\E X_\ell)/\sqrt{V_\ell}\dto\Normal$.
\end{theorem}
\begin{proof}
Fix $\epsilon>0$ and put $a_\ell=\epsilon\sqrt{V_\ell}$. Delete every
arriving batch whose cardinality exceeds $a_\ell$, or equivalently replace
it by an empty batch. Denote the retained discovery count by $X'_\ell$.
Under this coupling, $0\le X_\ell-X'_\ell\le R_\ell$, where $R_\ell$ is
the total cardinality of the removed arrivals. Set
$u_\ell=t_\ell\E[J_\ell^2;J_\ell>a_\ell]$.
The compound-Poisson second-moment identity gives
\[
 \E R_\ell^2=u_\ell+
       (t_\ell\E[J_\ell;J_\ell>a_\ell])^2.
\]
Since $J_\ell\le J_\ell^2/a_\ell$ on the deleted event,
\begin{equation}\label{eq:truncationL2}
 \frac{\E R_\ell^2}{V_\ell}
 \le\frac{u_\ell}{V_\ell}
   +\frac1{\epsilon^2}\left(\frac{u_\ell}{V_\ell}\right)^2\to0.
\end{equation}
Thus the centered difference has $L^2$ norm $o(\sqrt{V_\ell})$ and
$V'_\ell=\Var(X'_\ell)\sim V_\ell$. Eventually the retained model has
integer batch bound at most $\max(1,\lfloor a_\ell\rfloor)$. Applying
Theorem~\ref{thm:BE} gives a limiting upper bound $C\epsilon^{2/5}$ for its
Kolmogorov error. The $L^2$ coupling, normal continuity, and $V'_\ell\sim V_\ell$
transfer the same limiting bound to $X_\ell$; explicitly, one may first allow
a fixed smoothing tolerance and then let that tolerance decrease to zero.
Finally let $\epsilon\downarrow0$.
\end{proof}

\begin{corollary}[Fixed-law, fixed-sample extension]\label{cor:unboundedfixed}
Let $\mu$ be fixed with $\E|A_1|^2<\infty$, and suppose $V(n)\to\infty$,
equivalently $s_n^2\to\infty$. If for every $\epsilon>0$,
\[
 n\E[|A_1|^2;|A_1|>\epsilon\sqrt{V(n)}]=o(V(n)),
\]
then $(D_n-b_n)/s_n\dto\Normal$.
\end{corollary}
\begin{proof}
Apply Theorem~\ref{thm:unbounded} at $t=n$, then use
Proposition~\ref{prop:variance-equivalence}, Lemma~\ref{lem:window}, and the
coupling proof of Theorem~\ref{thm:main}. All these transfer steps require
only the stated moments and a fixed law. This does not claim a uniform
fixed-sample triangular-array theorem.
\end{proof}

\begin{example}[A finite second moment does not imply a CLT]\label{ex:unbounded}
There is a fixed law on disjoint finite blocks for which $\E|A_1|^2<\infty$
and $V(t)\to\infty$, but the variance-normalized discovery count has a
non-Gaussian two-point subsequential limit.
\end{example}
\begin{proof}
Take disjoint blocks $C_k$ with $s_k=|C_k|=2^{k^2}$ and probabilities
$\mu_k=c\,2^{-2k^2-k}$, $k\ge1$, where $c$ normalizes their sum.
Then $s_k^2\mu_k=c2^{-k}$, so $\E|A_1|^2<\infty$.
The independent block-discovery indicators $I_k(t)$ are Bernoulli with
parameters $1-e^{-t\mu_k}$, and
\[
 X(t)=\sum_k s_k I_k(t),\qquad
 V(t)=\sum_k s_k^2 e^{-t\mu_k}(1-e^{-t\mu_k}).
\]
Set $t_K=1/\mu_K$, $p=1-e^{-1}$, and $q=p(1-p)$.
The block $K$ contributes $q s_K^2$ to $V(t_K)$. For $k<K$, the total
variance is at most $K s_{K-1}^2e^{-\mu_{K-1}/\mu_K}=o(s_K^2)$.
Since $\max_{k>K}t_K\mu_k\to0$, the later blocks contribute
\[
 (1+o(1))t_K\sum_{k>K}s_k^2\mu_k
 =(1+o(1))t_Kc2^{-K}=(1+o(1))s_K^2.
\]
Thus $V(t_K)/s_K^2\to q+1$. On the other hand,
\[
 \Pp(\exists k>K:I_k(t_K)=1)\le t_K\sum_{k>K}\mu_k\to0,
 \qquad
 \frac{t_K\sum_{k>K}s_k\mu_k}{s_K}=O(2^{-2K})\to0.
\]
The centered later-block contribution divided by $s_K$ therefore vanishes
in probability, as does the earlier contribution in $L^2$. It follows that
\begin{equation}\label{eq:twopoint}
 \frac{X(t_K)-m(t_K)}{\sqrt{V(t_K)}}
 \dto \frac{Z-p}{\sqrt{q+1}},\qquad Z\sim\operatorname{Bernoulli}(p).
\end{equation}
The limiting variance is $q/(q+1)<1$: rare large blocks retain second moment
but disappear from the weak limit, so normalized squares are not uniformly
integrable. Finally, if $t_K\le t<t_{K+1}$, then $t\mu_k\le1$ for $k>K$.
Using $e^{-u}(1-e^{-u})\ge e^{-1}(1-e^{-1})u$ on $[0,1]$ gives
\[
 V(t)\ge e^{-1}(1-e^{-1})t\sum_{k>K}s_k^2\mu_k
 \ge e^{-1}(1-e^{-1})s_K^2\to\infty.
\]
Thus divergence holds for all $t\to\infty$, not only the subsequence.
Taking $n_K=\lfloor t_K\rfloor$ preserves the preceding calculations;
the general square-root-window coupling and variance equivalence transfer
the same non-Gaussian limit to the fixed-sample count.
\end{proof}

The example shows the need for additional control beyond a finite second
moment. It does not assert necessity of the sufficient condition
\eqref{eq:batchlindeberg}. That condition deliberately controls all incidence
mass in large batches, including already-discovered species.

\section{Actual-variance concentration and moderate deviations}\label{sec:MDP}

\begin{proposition}[Actual-variance Bernstein bound]\label{prop:actualmgf}
Let $M=X(t)-m(t)$, $V=V(t)$, and
$\psi_B(u)=(e^{Bu}-1-Bu)/B^2$ for $u\ge0$. For every $q>1$ and real $\theta$,
\begin{equation}\label{eq:holdermgf}
 \log\E e^{\theta M}
 \le\frac{q-1}{q}\frac V{L_B}
 \left[\exp\left\{\frac{L_B\psi_B(q|\theta|)}{q-1}\right\}-1\right].
\end{equation}
In particular,
\begin{equation}\label{eq:smallmgf}
 \log\E e^{\theta M}\le
 \frac{V\theta^2}{2(1-2B|\theta|)},\qquad |\theta|<\frac1{2B},
\end{equation}
and, for every $x>0$,
\begin{equation}\label{eq:actualconcentration}
 \Pp(|X(t)-m(t)|\ge x)
 \le2\exp\left\{-\frac{x^2}{2(V+2Bx)}\right\}.
\end{equation}
Separately for either sign and for every $\epsilon>0$,
\begin{equation}\label{eq:neargaussian}
 \Pp(\pm M\ge x)\le
 e^{-x^2/(2(1+\epsilon)V)}
 \quad\text{if }0<x\le\frac{\epsilon V}{2B}.
\end{equation}
\end{proposition}
\begin{proof}
The compensated-Poisson representation gives a mean-one exponential martingale
\begin{equation}\label{eq:exp-mart}
 Z_t(\theta)=e^{\theta M-K_t(\theta)},\qquad
 K_t(\theta)=\int_0^t\sum_C\mu(C)
       (e^{\theta H_C(s)}-1-\theta H_C(s))\,ds.
\end{equation}
For example, its stochastic differential is obtained by the jump product rule;
the compensated integral has zero drift. It is a true martingale: bounded
integrands, finite total arrival rate, and finite horizon dominate its stopped
versions by an integrable exponential of $N(t)+t$. Moreover,
$0\le K_t(\theta)\le\psi_B(|\theta|)Q_t$ by the exponential series and
$0\le H_C\le B$.
H\"older with exponents $q$ and $q/(q-1)$ therefore gives
\[
 \E e^{\theta M}
 \le\left(\E e^{q\theta M-\psi_B(q|\theta|)Q_t}\right)^{1/q}
      \left(\E e^{\psi_B(q|\theta|)Q_t/(q-1)}\right)^{(q-1)/q}
 \le\left(\E e^{\psi_B(q|\theta|)Q_t/(q-1)}\right)^{(q-1)/q}.
\]
The uncentered version of \eqref{eq:bracketmgf} proves \eqref{eq:holdermgf}.

To prove \eqref{eq:smallmgf}, the case $\theta=0$ is immediate. Otherwise
set $a=B|\theta|\in(0,1/2)$ and $q=1+a$. The function
$(e^{Lz}-1)/L$ increases in $L>0$ for $z\ge0$, so we may replace $L_B$
by $2B^2$ in \eqref{eq:holdermgf}. The elementary series bounds
\[
 \psi(u)\le\frac{u^2}{2(1-u/3)}\quad(0\le u<3),\qquad
 e^z-1\le\frac z{1-z/2}\quad(0\le z<2)
\]
then yield
\[
 \log\E e^{\theta M}
 \le\frac{V\theta^2}{2}
       \frac q{1-qa/3-q^2a/2}.
\]
All denominators are positive: for $q=1+a$,
\[
 1-\frac{qa}{3}-\frac{q^2a}{2}-q(1-2a)
 =\frac a6+\frac{2a^2}{3}-\frac{a^3}{2}\ge0.
\]
In particular the argument of the second series bound is less than two,
and the displayed estimate implies \eqref{eq:smallmgf}.
Choose $|\theta|=x/(V+2Bx)$ in the Chernoff bound for each sign to obtain
\eqref{eq:actualconcentration}. If $2Bx\le\epsilon V$, this also proves
\eqref{eq:neargaussian} without the factor two for each individual tail.
\end{proof}

\begin{remark}[The near-Gaussian range]\label{rem:gaussianrange}
The quadratic variance coefficient in \eqref{eq:smallmgf} is the actual $V$.
Its first correction has order $B|\theta|$, not $B^2\theta^2$.
A universal $1+O(\epsilon)$ relative moment bound throughout
$|\theta|B\asymp\sqrt\epsilon$ would be false: for
$Z\sim\operatorname{Bernoulli}(p)$, $0<p<1/2$,
\[
 \frac{\log\E e^{\theta(Z-p)}}{p(1-p)\theta^2/2}
 =1+\frac{1-2p}{3}\theta+O(\theta^2).
\]
Such Bernoulli sums occur already in Poissonized singleton discovery.
The range $x\le\epsilon V/(2B)$ in \eqref{eq:neargaussian} is the one
proved here; no $\sqrt\epsilon$ replacement is claimed.
\end{remark}

The direct Bennett estimate \eqref{eq:discoverybennett} and
\eqref{eq:actualconcentration} are complementary. The latter has the exact
quadratic variance coefficient; the former can be stronger at larger deviations
and can be refined by the fractional allocation in
Proposition~\ref{prop:fractionalbennett}. Neither is substituted for a
fixed-sample actual-variance bound.

\begin{theorem}[Gaussian moderate-deviation principle]\label{thm:MDP}
Consider any sequence of Poissonized bounded-batch models, with counts $X_k$,
variances $V_k>0$, and batch bounds $B_k$. Suppose
\[
 a_k\longrightarrow\infty,\qquad
 \frac{a_kB_k}{\sqrt{V_k}}\longrightarrow0.
\]
Then
\[
 \frac{X_k-\E X_k}{a_k\sqrt{V_k}}
\]
satisfies a large-deviation principle with speed $a_k^2$ and good rate function
$I(x)=x^2/2$. In particular, for each fixed $x>0$,
\begin{equation}\label{eq:MDPtail}
 \frac1{a_k^2}\log\Pp(X_k-\E X_k\ge xa_k\sqrt{V_k})
 \longrightarrow-\frac{x^2}{2},
\end{equation}
and the analogous lower-tail limit holds.
\end{theorem}
\begin{proof}
Write $M_k=X_k-\E X_k$, let $Q_k$ be its terminal bracket, and fix
$\lambda\in\R$. Set $\theta_k=\lambda a_k/\sqrt{V_k}$.
For each fixed $\delta\in(0,1)$ let
$G_k=\{|Q_k/V_k-1|\le\delta\}$. By
\eqref{eq:Qupper}--\eqref{eq:Qlower} and $L_{B_k}\le2B_k^2$,
\begin{equation}\label{eq:superexpbracket}
 \Pp(G_k^c)\le2\exp\{-c_\delta V_k/B_k^2\},
 \qquad V_k/B_k^2\gg a_k^2.
\end{equation}
Uniform Taylor expansion for $|\theta_k|B_k=o(1)$ gives
\begin{equation}\label{eq:Kexpansion}
 K_k(\theta_k)=\frac{\theta_k^2Q_k}{2}
           (1+O(|\theta_k|B_k)).
\end{equation}
The relative error bound follows termwise from
$e^z-1-z=z^2(1+O(|z|))/2$ for real $z$ near zero; it is uniform over
all batch types and histories.

One must also control the exceptional event under the exponential change of
measure. By \eqref{eq:smallmgf}, for large $k$,
\[
 \E e^{2\theta_kM_k}\le e^{4\lambda^2a_k^2}.
\]
Since $K_k(\theta_k)\ge0$, Cauchy--Schwarz gives both
\begin{align}
 \E[e^{\theta_kM_k};G_k^c]
 &\le\exp\{2\lambda^2a_k^2\}\Pp(G_k^c)^{1/2},\label{eq:bad-original}\\
 \E[Z_k(\theta_k);G_k^c]
 &\le\exp\{2\lambda^2a_k^2\}\Pp(G_k^c)^{1/2}.
 \label{eq:bad-tilt}
\end{align}
These quantities are superexponentially small at speed $a_k^2$.
In particular, $\E[Z_k(\theta_k);G_k]=1-o(1)$.
On $G_k$, \eqref{eq:Kexpansion} bounds $K_k(\theta_k)$ between
\[
 \frac{\lambda^2a_k^2}{2}(1-\delta)(1-o(1))
 \quad\text{and}\quad
 \frac{\lambda^2a_k^2}{2}(1+\delta)(1+o(1)).
\]
Use $e^{\theta_kM_k}=Z_k(\theta_k)e^{K_k(\theta_k)}$, together with
\eqref{eq:bad-original}--\eqref{eq:bad-tilt}, to bound the moment generating
function from above and below. After letting $\delta\downarrow0$,
\begin{equation}\label{eq:limitingcgf}
 \lim_{k\to\infty}\frac1{a_k^2}
 \log\E\exp\left\{\frac{\lambda a_kM_k}{\sqrt{V_k}}\right\}
 =\frac{\lambda^2}{2}.
\end{equation}
The G\"artner--Ellis theorem~\cite[Theorem~2.3.6]{DZ}, whose limiting cumulant generating
function here is finite and differentiable on all of $\R$, proves the stated
principle. The half-line limits follow because the interior and closure have
the same infimum of $I$ at a positive threshold.
\end{proof}

This proof is consistent with the general bounded-jump martingale
moderate-deviation theory~\cite[Corollary~1(a)]{Dembo}; the additional point supplied by the
model is the exponential bracket estimate. Ordinary convergence
$Q_k/V_k\to1$ in probability alone would not justify the change-of-measure
step. Moreover, \eqref{eq:MDPtail} is a \emph{logarithmic} tail asymptotic.
It does not assert that the ratio to the corresponding Gaussian tail tends to
one. No fixed-sample moderate-deviation principle is claimed here.

\section{Multiple-coverage counts}\label{sec:threshold}

For $r\in\N$, let
\[
 K_r(t)=\sum_i\ind_{\{N_i(t)\ge r\}},\qquad
 q_{i,r}(t)=\Pp(\Poi(t\lambda_i)\ge r),\qquad
 v_r(t)=\sum_iq_{i,r}(t)(1-q_{i,r}(t)).
\]
Here $N_i(t)$ counts batches containing species $i$, and $K_1(t)=X(t)$.

\begin{proposition}[Threshold variance and Bennett bounds]\label{prop:thresholdvar}
With $V_r(t)=\Var(K_r(t))$,
\begin{equation}\label{eq:thresholdvar}
 v_r(t)\le V_r(t)\le Bv_r(t).
\end{equation}
The concentration inequalities \eqref{eq:discoverybennett}--\eqref{eq:discoverybernstein}
hold with $X,m,v_0$ replaced by $K_r,\E K_r,v_r$.
\end{proposition}
\begin{proof}
The indicators are increasing functions of independent Poisson batch-type
counts, and hence have nonnegative covariances by product-measure association
\cite{EPW}. Alternatively, this instance follows by induction over independent
coordinates from
$\Cov(g(U),h(U))=\tfrac12\E[(g(U)-g(U'))(h(U)-h(U'))]\ge0$
for increasing $g,h$ and an independent copy $U'$; bounded conditional
expectations give the countable-coordinate passage.
This gives the lower bound. The read-$B$ variance and exponential inequalities
give the remaining assertions. Countable approximation is justified by
$K_r(t)\le X(t)\le BN(t)$.
\end{proof}

\begin{theorem}[Poissonized multiple-coverage approximation]\label{thm:thresholdCLT}
For $r\ge1$, define
\begin{equation}\label{eq:dr}
 d_1=1,\qquad d_r=\min\{2^{2r-1},e^2r\}\quad(r\ge2).
\end{equation}
For every $t$ with $V_r(t)>0$ and every bounded-batch law,
\begin{equation}\label{eq:thresholdBE}
 \dK\left(\frac{K_r(t)-\E K_r(t)}{\sqrt{V_r(t)}},\Normal\right)
 \le\min\left\{1,C\left(\frac{rB^2}{V_r(t)}\right)^{1/5}\right\},
\end{equation}
where $C$ is absolute. In particular, along any sequence of laws, times,
batch bounds, and thresholds such that $rB^2/V_r(t)\to0$,
\begin{equation}\label{eq:thresholdCLT}
 \frac{K_r(t)-\E K_r(t)}{\sqrt{V_r(t)}}\dto\Normal.
\end{equation}
Thus the assertion includes growing thresholds, not only fixed $r$.
\end{theorem}
\begin{proof}
For a fixed terminal time $t$, consider the terminal Doob martingale of
$K_r(t)$. Just before time $s$, the increase in the conditional expectation of
the contribution of species $i$, if that species is included in an arriving
batch, is
\begin{equation}\label{eq:threshold-h}
 h_{i,r}(s)=\ind_{\{N_i(s-)<r\}}
 \Pp\{\Poi(\lambda_i(t-s))=r-1-N_i(s-)\}.
\end{equation}
This follows by subtracting the two Poisson upper-tail probabilities before
and after increasing the current count by one. Thus the batch jump is
$H_{C,r}(s)=\sum_{i\in C}h_{i,r}(s)$ and lies in $[0,B]$.
The bracket is
\begin{equation}\label{eq:thresholdQ}
 Q_{r,t}=\sum_{i,j}\lambda_{ij}\int_0^t h_{i,r}(s)h_{j,r}(s)\,ds,
 \qquad \E Q_{r,t}=V_r(t).
\end{equation}

For $0\le k\le r-1$ and $x\ge0$,
\[
 e^{-x}\frac{x^k}{k!}
 =2^k e^{-x/2}\left[e^{-x/2}\frac{(x/2)^k}{k!}\right]
 \le2^{r-1}e^{-x/2}.
\]
For $r\ge2$, also put $\vartheta=1-1/r$. Since a Poisson point probability
is at most one,
\[
 e^{-x}\frac{x^k}{k!}
 =\vartheta^{-k}e^{-x/r}
       \left[e^{-\vartheta x}\frac{(\vartheta x)^k}{k!}\right]
 \le\left(\frac r{r-1}\right)^{r-1}e^{-x/r}
 \le e\,e^{-x/r}.
\]
Integrating the products of either envelope, and using the exact discovery
bound when $r=1$, gives
\begin{equation}\label{eq:threshold-envelope}
 0\le\int_0^t h_{i,r}(s)h_{j,r}(s)\,ds
 \le\frac{d_r}{\lambda_i+\lambda_j}.
\end{equation}
Normalize the integral in \eqref{eq:threshold-envelope} to a random variable
$F_{ij}^{(r)}\in[0,1]$. It is a function of the \emph{entire incident Poisson
paths}, not just their first arrivals. We obtain the positive decomposition
\[
 Q_{r,t}=\sum_{i,j}d_r a_{ij}F_{ij}^{(r)}.
\]
Each batch-path coordinate has coefficient load at most $d_rL_B$. Therefore,
\begin{align}
 \Var(Q_{r,t})&\le d_rL_BV_r(t)\le2e^2rB^2V_r(t),\label{eq:thresholdbracketvar}\\
 \log\E e^{\theta(Q_{r,t}-V_r(t))}
 &\le\frac{V_r(t)}{d_rL_B}\psi(d_rL_B\theta).
 \label{eq:thresholdbracketmgf}
\end{align}
Our choice $d_1=1$ retains the exact discovery constant.
The bounded jumps and the fourth-moment argument of
Theorem~\ref{thm:BE} now prove \eqref{eq:thresholdBE} and hence
\eqref{eq:thresholdCLT}. The deterministic-grid justification applies because
these integrands are bounded and piecewise smooth between the finitely many
batch arrivals, with the explicit bracket-density derivative bound in
Appendix~\ref{app:mesh}. All fourth-moment constants are universal;
$r$ enters only through $d_r\le e^2r$.
\end{proof}

\begin{corollary}[Further threshold fluctuations]\label{cor:thresholdMDP}
The logarithmic bound also extends in the form
$\dK\le\min\{1,C[\eta_r\log(e+\eta_r^{-1})]^{1/4}\}$, where
$\eta_r=rB^2/V_r(t)$ and $C$ is universal.
For each fixed $r$, the higher-moment normal-approximation estimates and the
Poissonized moderate-deviation principle extend to $K_r$, with constants
allowed to depend on $r$. In particular, the moderate-deviation range remains
$a\to\infty$, $aB/\sqrt{V_r}\to0$.
\end{corollary}
\begin{proof}
For the logarithmic bound, repeat Theorem~\ref{thm:BElog}: the bracket tails
have exponent at least $c\delta^2/\eta_r$, while the deterministic-bracket
completion error is at most $C(B^2/V_r)^{2/7}\le C\eta_r^{2/7}$.
For the other assertions, use \eqref{eq:thresholdbracketmgf} in the proofs of
Corollary~\ref{cor:brackettails} and Theorem~\ref{thm:MDP}. The overlap constant
is at most $2e^2rB^2$, a fixed multiple of $B^2$ for fixed $r$, and the martingale
jump bound remains $B$. In the small-moment bound the constants depend on $r$;
this does not affect the exponential-negligibility comparison in
\eqref{eq:superexpbracket}.
\end{proof}

\begin{remark}[What is not obtained by subtraction]
The number of species observed exactly $r$ times is
$K_r(t)-K_{r+1}(t)$. Separate scalar CLTs for these two variables do not imply
its CLT. Their fluctuations can cancel, and signed cross-bracket terms are not
covered by the preceding positive decomposition. The read-$B$ variance upper
bound and exponential-moment comparison do apply to exact-frequency counts,
but their universal variance-normalized Gaussian limits are not proved here.
Nor is a fixed-sample version of Theorem~\ref{thm:thresholdCLT} asserted.
\end{remark}

\section{Weighted and conditional discoveries; temporal covariance}\label{sec:extensions}

\subsection{Nonnegative weights and bounded total contribution}

\begin{theorem}[Weighted discovery]\label{thm:weighted}
Assign weights $w_i\ge0$ and assume that every positive-probability batch obeys
\begin{equation}\label{eq:weightedbatch}
 \sum_{i\in C}w_i\le W<\infty.
\end{equation}
Batch cardinalities need not be uniformly bounded, although batches remain
finite. Let
$X_w(t)=\sum_iw_i\ind_{\{T_i\le t\}}$ and
$V_w(t)=\Var(X_w(t))$. Its terminal martingale has jumps at most $W$ and bracket
$Q_{w,t}$ satisfying
\begin{align}
 \Var(Q_{w,t})&\le2W^2V_w(t),\label{eq:weightedQvar}\\
 \log\E e^{\theta(Q_{w,t}-V_w(t))}
 &\le\frac{V_w(t)}{2W^2}\psi(2W^2\theta).
 \label{eq:weightedQmgf}
\end{align}
For $W>0$, the Poissonized normal-approximation, actual-variance concentration,
and moderate-deviation results above hold with $B$ replaced by $W$. They also
hold for sequences with varying weights and laws subject to the corresponding
conditions on $W^2/V_w$.
\end{theorem}
\begin{proof}
The bound $X_w(t)\le WN(t)$ gives every required moment. The jump integrand is
$\sum_{i\in C}w_i e^{-\lambda_i(t-s)}Y_i(s-)$, and
\[
 Q_{w,t}=\sum_{i,j}w_iw_j a_{ij}F_{ij}.
\]
The incidence identity becomes
\[
 \sum_jw_j\lambda_{ij}
 =\E\left[\ind_{\{i\in A_1\}}\sum_{j\in A_1}w_j\right]
 \le W\lambda_i.
\]
Hence the coefficient load of coordinate $C$ is at most
\[
 2\sum_{i\in C}\frac{w_i}{\lambda_i}\sum_jw_j\lambda_{ij}
 \le2W\sum_{i\in C}w_i\le2W^2.
\]
Apply the weighted-overlap inequalities. The remaining proofs use only that
load bound and the jump bound, so they transfer verbatim. The case $W=0$ is
constant and needs no normalization.
\end{proof}

Nonnegativity is essential to this particular argument: it ensures a positive
bracket decomposition with expectation $V_w$. The signed-weight variance bound
\eqref{eq:signedweightedvariance} does not, by itself, provide the same
actual-variance CLT for arbitrary signed weights.

\subsection{Future discoveries conditional on the observation history}
Let
$\Delta_{t,u}=X(t+u)-X(t)$ for $u>0$, and define
\[
 m_{t,u}=\E[\Delta_{t,u}\mid\mathcal F_t],\qquad
 V_{t,u}=\Var(\Delta_{t,u}\mid\mathcal F_t).
\]

\begin{corollary}[Conditional oracle normal approximation]\label{cor:conditional}
Almost surely on $\{V_{t,u}>0\}$,
\begin{equation}\label{eq:conditionalBE}
 \sup_x\left|\Pp\left(
 \left.\frac{\Delta_{t,u}-m_{t,u}}{\sqrt{V_{t,u}}}\le x
 \right|\mathcal F_t\right)-\Phi(x)\right|
 \le\mathcal R(B^2/V_{t,u}).
\end{equation}
\end{corollary}
\begin{proof}
Condition on the past, and intersect each future batch with the set of species
unseen at time $t$. Independent increments make the future a new Poissonized
bounded-batch discovery experiment. Apply Theorem~\ref{thm:BElog} to its
conditional law. The constant is universal and independent of the past.
\end{proof}

For example,
$m_{t,u}=\sum_{i:T_i>t}(1-e^{-u\lambda_i})$, and $V_{t,u}$ is obtained from
\eqref{eq:exactvariance} by restricting both indices to the unseen set and
replacing $t$ by $u$. These are oracle quantities involving the unknown law.
The corollary does not construct estimators for them.

\subsection{Covariance across time and an increment bound}

\begin{proposition}[Temporal covariance]\label{prop:temporal}
For $s,t\ge0$,
\begin{equation}\label{eq:temporal-cov}
 \Cov(X(s),X(t))
 =\sum_{i,j}e^{-s\lambda_i-t\lambda_j}
       (e^{\min(s,t)\lambda_{ij}}-1).
\end{equation}
For $0\le s\le t$,
\begin{equation}\label{eq:incrementvar}
 \Var(X(t)-X(s))\le B\sum_i p_i(1-p_i)
 \le B(m(t)-m(s)),\qquad p_i=e^{-s\lambda_i}-e^{-t\lambda_i}.
\end{equation}
The read-$B$ exponential comparison also applies to these increments.
\end{proposition}
\begin{proof}
For $s\le t$, the joint absence probability is
\[
 \Pp(T_i>s,T_j>t)=e^{-s\lambda_i-t\lambda_j+s\lambda_{ij}}.
\]
Subtract the product of the marginal absence probabilities and sum; symmetry
covers the other order. For the increment bound, apply
Corollary~\ref{cor:readB} to $\ind_{\{s<T_i\le t\}}$, whose support is still
the set of incident batch clocks. No association claim is needed for these
nonmonotone interval indicators.
\end{proof}

\section{Scope and remaining questions}\label{sec:scope}

\paragraph{Local limits and arithmetic structure.}
If the batch types are disjoint pairs, $X(t)$ and $D_n$ are always even.
Thus variance divergence and a CLT do not imply a span-one local limit theorem.
A lattice-corrected theorem needs additional arithmetic and Fourier analysis,
including control of possible approximate periodicity.

\paragraph{Relative Gaussian tails.}
Theorem~\ref{thm:MDP} identifies logarithmic probabilities on a growing scale.
It does not show
$\Pp(X-m\ge x\sqrt V)/(1-\Phi(x))\to1$. A Cram\'er-type relative approximation,
and its largest valid range, remain separate questions. An optimal
$O(B/\sqrt V)$ Kolmogorov bound is also not established by the generic
martingale estimate used here.

\paragraph{Fixed-sample refinements.}
The fixed-law CLT is proved with the actual fixed-sample variance. Nevertheless,
a variance-only uniform fixed-sample Berry--Esseen bound, actual-variance
fixed-sample concentration, and fixed-sample moderate deviations are not proved
by the present transfer argument. A vanishing-probability coupling error is
not automatically superexponentially negligible.

\paragraph{Occupancy spectra and process limits.}
The read-$B$ calculus controls species-wise bounded occupancy functionals,
including exact-frequency counts and bounded-coefficient truncated
Good--Toulmin statistics~\cite{GoodToulmin}. It does not automatically give their joint Gaussian
limits. Likewise, \eqref{eq:temporal-cov} and \eqref{eq:incrementvar} are useful
inputs for a functional CLT, but covariance convergence alone does not prove
finite-dimensional Gaussianity and tightness. Signed linear combinations and
variance cancellation must be handled explicitly.

\paragraph{Statistical calibration and random environments.}
Unknown-law prediction intervals require joint fluctuation theorems, bias
control, and consistent variance estimation, not merely the oracle normal
approximation in Corollary~\ref{cor:conditional}. Conditioning on a random law
also does not generally give an unconditional Gaussian limit under deterministic
centering and scaling; random asymptotic means or variances may survive. Finally,
the finite-second-moment variance expansion in
Theorem~\ref{thm:correction} is not a CLT for all unbounded batch laws.
The weighted theorem and the Lindeberg theorem supply two precise extensions
beyond bounded cardinality. Example~\ref{ex:unbounded} rules out an
unrestricted finite-second-moment conclusion.

\appendix
\section{Deterministic-grid justification of the martingale bound}\label{app:mesh}

For completeness, we justify the use of \eqref{eq:Haeusler} without assuming a
continuous-time formulation. Fix one model and one terminal time $t$, and let
$M$ be its terminal martingale. Write
\[
 M_s=\sum_{u\le s}\Delta M_u-\int_0^s c_u\,du,
 \qquad Q_s=\int_0^s q_u\,du.
\]
For discovery, $0\le c_u\le B$, $0\le q_u\le B^2$, and there are only finitely
many jumps on $[0,t]$, almost surely. Between batch arrivals $q_u$ is Lipschitz,
with a bound depending on $B$: differentiating the exponential integrands gives
$|q_u'|\le2B^2$. The same constant works for every threshold $r$. Indeed, with
$p_k(x)=e^{-x}x^k/k!$ and $p_{-1}=0$, one has
$p_k'(x)=p_{k-1}(x)-p_k(x)$ and $|p_k'(x)|\le1$.
Between arrivals, $|h_{i,r}'(s)|\le\lambda_i\le1$, hence
$|H_{C,r}'(s)|\le B$ and $|q_s'|\le2B^2$.
Weighted discovery is identical with $B$ replaced by $W$, since
$\sum_{i\in C}w_i\lambda_i\le W$.

Let $\pi$ be a deterministic partition
$0=t_0<t_1<\cdots<t_m=t$ with mesh $|\pi|\to0$, and set
\[
 \zeta_j=M_{t_j}-M_{t_{j-1}},\qquad
 Q^\pi=\sum_j\E[\zeta_j^2\mid\mathcal F_{t_{j-1}}].
\]
These are genuine martingale differences, but they are \emph{not} asserted to
be bounded by $B$; a grid interval can contain several arrivals.
The conditional isometry gives
\[
 Q^\pi=\sum_j\int_{t_{j-1}}^{t_j}
                    \E[q_s\mid\mathcal F_{t_{j-1}}]\,ds.
\]
For $p\ge2$, boundedness, the Lipschitz behavior between arrivals, and the
probability at most $|\pi|$ of any arrival in one interval imply
\[
 \|q_s-q_{t_{j-1}}\|_p
 \le C_B(|\pi|+|\pi|^{1/p}),\qquad t_{j-1}\le s\le t_j.
\]
Conditional expectation is an $L^p$ contraction, so integration gives
\begin{equation}\label{eq:gridQlimit}
 \|Q^\pi-Q_t\|_p\le2tC_B(|\pi|+|\pi|^{1/p})\longrightarrow0.
\end{equation}
Values at deterministic endpoints cause no issue, since their jump
probabilities are zero.

Pathwise, a finite-jump process with a Lipschitz continuous drift satisfies,
for $p>1$,
\[
 \sum_j|M_{t_j}-M_{t_{j-1}}|^{2p}
 \longrightarrow\sum_{s\le t}|\Delta M_s|^{2p}.
\]
Indeed, sufficiently fine intervals separate the finitely many jumps; the
continuous contribution to each such increment tends to zero, and the sum of
$2p$-th powers of increments on jump-free intervals tends to zero. The whole
sum is bounded by the $2p$-th power of total variation,
\[
 \sum_j|M_{t_j}-M_{t_{j-1}}|^{2p}
 \le B^{2p}(N(t)+t)^{2p},
\]
which is integrable. Hence
\begin{equation}\label{eq:gridjumplimit}
 \sum_j\E|M_{t_j}-M_{t_{j-1}}|^{2p}
 \longrightarrow\E\sum_{s\le t}|\Delta M_s|^{2p}.
\end{equation}

Apply \eqref{eq:Haeusler} to the grid differences divided by $\sqrt V$.
Their sum is exactly $M_t/\sqrt V$ for every partition, and their total
unconditional variance is exactly one. Taking the mesh limit using
\eqref{eq:gridQlimit}--\eqref{eq:gridjumplimit} proves
\begin{equation}\label{eq:continuousHaeusler}
 \dK\left(\frac{M_t}{\sqrt V},\Normal\right)
 \le C_p\left\{
 \E\left|\frac{Q_t}{V}-1\right|^p
 +\frac{\E\sum_{s\le t}|\Delta M_s|^{2p}}{V^p}
 \right\}^{1/(2p+1)}.
\end{equation}
The mesh limit is taken separately for each model. Uniformity of the final
normal-approximation bound comes from the moment estimates in the main text,
not from a uniform-in-$t$ discretization estimate.

\subsection*{Stopping and Gaussian completion}
For the stopped martingale $M_{s\wedge\tau}$ in
Lemma~\ref{lem:completion}, the bracket density is
$q_s\ind_{\{s\le\tau\}}$, up to irrelevant point values.
It remains bounded and has only finitely many discontinuities along each
sample path (batch arrivals and the possible additional stopping time).
The explicit probability bound used above for an arrival need not hold for
that stopping time. Instead, write $\ell_\pi(s)$ for the left grid endpoint.
Write $q_s^{\mathrm{stop}}=q_s\ind_{\{s\le\tau\}}$.
For almost every $(s,\omega)$ with respect to Lebesgue measure times
probability, the left-grid values converge to $q_s^{\mathrm{stop}}$.
Boundedness, dominated convergence on $[0,t]\times\Omega$, and H\"older
in $s$ therefore give
\[
 \int_0^t\|q_s^{\mathrm{stop}}-
                   q_{\ell_\pi(s)}^{\mathrm{stop}}\|_p\,ds\to0.
\]
The conditional-expectation contraction then yields
$\|Q^\pi-Q_\tau\|_p\to0$. The stopped path still has total variation at most
$B(N(t)+t)$, so the summed high-moment limit \eqref{eq:gridjumplimit}
also holds with stopping. This is exactly the passage used before appending
the Gaussian differences in Lemma~\ref{lem:completion}.

\section{Revision audit and reproducibility}\label{app:audit}

This revision contains proofs of the retained theorem statements and has
undergone an internal proof and reference audit. This is not formal verification
or independent peer review. The literature comparison is limited to the cited
sources; no exhaustive priority search is claimed.

\renewcommand{\arraystretch}{1.14}
\begin{longtable}{@{}>{\raggedright\arraybackslash}p{0.25\textwidth}>{\raggedright\arraybackslash}p{0.69\textwidth}@{}}
\toprule
Issue & Treatment in this revision\\
\midrule
\endfirsthead
\toprule
Issue & Treatment in this revision\\
\midrule
\endhead
\bottomrule
\endfoot
Product-space limits & Lemma~\ref{lem:finner} specifies the finite-family,
finite-coordinate, and truncation limits. Centered $L^2$ series and nonnegative
sums are distinguished.\\
Species refinements & Proposition~\ref{prop:speciesvar} allows signed weights
and conditional expected batch sizes. Proposition~\ref{prop:fractionalbennett}
uses fractional support allocations, not a frequency-sensitive interpretation
of the maximum supported batch size.\\
Bracket load & Ordered pairs include the diagonal $a_{ii}=1/2$.
Proposition~\ref{prop:loadsharp} proves that $L_B$ is optimal for this
worst-case coefficient-load calculation, not for every possible bracket bound.\\
Normal approximation & Haeusler's bound as recorded by Mourrat is used without
assuming bounded grid increments. The logarithmic refinement is proved using
stopping, Gaussian completion, and a separate smoothing lemma, with fixed
moment order $p=3$. The stopped-grid limit is justified in Appendix~\ref{app:mesh}.\\
Unbounded batches & The Lindeberg truncation controls both terms in the
compound-Poisson second moment. Theorem~\ref{thm:unbounded} is Poissonized and
allows changing laws. Corollary~\ref{cor:unboundedfixed} is a fixed-law transfer.
Example~\ref{ex:unbounded} excludes an unrestricted finite-second-moment CLT.\\
Variance correction & The absolute $o(1)$ remainder retains diagonal terms.
Variance equivalence uses only a finite second moment; the window lemma uses
summability of the marginal rates, not a deterministic batch-size bound.\\
Actual-variance tails & Proposition~\ref{prop:actualmgf} gives the Bernstein
denominator $2(V+2Bx)$. The near-Gaussian range is linear in $\epsilon$,
not in $\sqrt\epsilon$. The rejected larger range is explained in
Remark~\ref{rem:gaussianrange}.\\
Thresholds & The bracket envelope has constant
$d_1=1$ and $d_r=\min(2^{2r-1},e^2r)$ for $r\ge2$.
Growing thresholds are permitted under $rB^2/V_r\to0$.
Full incident Poisson paths, rather than first arrivals, are the independent
coordinates.\\
Moderate deviations & The exceptional bracket event is controlled under
exponential tilting. The result is Poissonized and logarithmic; relative
Gaussian tails and fixed-sample moderate deviations are not claimed.\\
Attribution and scope & The independent-clock construction, generalized
H\"older/read-$B$ inequalities, association, infinite-urn derivative lemma,
self-bounding bounds, and generic martingale estimates are credited.
Unknown-law calibration and joint or functional Gaussian limits remain open
within the present argument.\\
\end{longtable}

\subsection*{Reproducible checks}
The accompanying \texttt{verify\_finite\_models.py} integrates finite-state
Kolmogorov and Feynman--Kac equations. It checks 40 bracket cases, 25 fixed-sample
covariance cases, and 80 centered-indicator moment bounds. All passed on rerun;
the maximum absolute discrepancy in $\E Q=\Var(K_r)$ was approximately
$4.49\times10^{-11}$.

The additional \texttt{verify\_revision.py} enumerates finite batch-presence
states and checks 108 signed-weight variance inequalities, 144 fractional
moment bounds, and 288 actual-variance moment bounds. It also verifies the
H\"older denominator identity symbolically, tests its parameter inequalities,
checks the explicit hub loads and threshold envelopes, and evaluates the
counterexample along nine parameter values. The JSON reports record the
parameters and outputs. These are deterministic numerical and symbolic sanity
checks, not proofs of the asymptotic or countable-space statements. The proofs
in the manuscript do not depend on their numerical outcomes.

A separate \texttt{reference\_audit.md} records the precise external results
used, source versions, and the distinction between direct theorem checks and
bibliographic or contextual checks. The source is self-contained: all
bibliographic entries are included below.

\clearpage
\phantomsection
\addcontentsline{toc}{section}{References}
\begin{thebibliography}{99}

\bibitem{BLM2000}
S. Boucheron, G. Lugosi, and P. Massart.
A sharp concentration inequality with applications.
\emph{Random Structures \& Algorithms} \textbf{16}(3) (2000), 277--292.
\href{https://doi.org/10.1002/(SICI)1098-2418(200005)16:3<277::AID-RSA4>3.0.CO;2-1}{Publisher DOI};
\href{https://econ-papers.upf.edu/papers/376.pdf}{author preprint}.

\bibitem{BLM}
S. Boucheron, G. Lugosi, and P. Massart.
On concentration of self-bounding functions.
\emph{Electronic Journal of Probability} \textbf{14} (2009), 1884--1899.
\href{https://doi.org/10.1214/EJP.v14-690}{doi:10.1214/EJP.v14-690}.

\bibitem{Dembo}
A. Dembo.
Moderate deviations for martingales with bounded jumps.
\emph{Electronic Communications in Probability} \textbf{1} (1996), 11--17.
\href{https://doi.org/10.1214/ECP.v1-973}{doi:10.1214/ECP.v1-973}.

\bibitem{DZ}
A. Dembo and O. Zeitouni.
\emph{Large Deviations Techniques and Applications}, 2nd ed.,
corrected printing of the 1998 edition.
Stochastic Modelling and Applied Probability, vol.~38. Springer, 2010.
\href{https://doi.org/10.1007/978-3-642-03311-7}{doi:10.1007/978-3-642-03311-7}.

\bibitem{Dutko}
M. Dutko.
Central limit theorems for infinite urn models.
\emph{Annals of Probability} \textbf{17}(3) (1989), 1255--1263.
\href{https://doi.org/10.1214/aop/1176991268}{doi:10.1214/aop/1176991268}.

\bibitem{ErikssonGraphs}
E. Eriksson.
Edge exchangeable graphs: connectedness, Gaussianity and completeness.
\href{https://arxiv.org/abs/2501.09511v2}{arXiv:2501.09511v2}, May 6, 2026.

\bibitem{ErikssonSpecies}
E. Eriksson.
The unseen species problem revisited.
\href{https://arxiv.org/abs/2602.08769v4}{arXiv:2602.08769v4}, September 9, 2026.

\bibitem{EPW}
J. D. Esary, F. Proschan, and D. W. Walkup.
Association of random variables, with applications.
\emph{Annals of Mathematical Statistics} \textbf{38}(5) (1967), 1466--1474.
\href{https://doi.org/10.1214/aoms/1177698701}{doi:10.1214/aoms/1177698701}.

\bibitem{Finner}
H. Finner.
A generalization of H\"older's inequality and some probability inequalities.
\emph{Annals of Probability} \textbf{20}(4) (1992), 1893--1901.
\href{https://doi.org/10.1214/aop/1176989534}{doi:10.1214/aop/1176989534}.

\bibitem{GLSS}
D. Gavinsky, S. Lovett, M. Saks, and S. Srinivasan.
A tail bound for read-$k$ families of functions.
\emph{Random Structures \& Algorithms} \textbf{47}(1) (2015), 99--108.
\href{https://doi.org/10.1002/rsa.20532}{doi:10.1002/rsa.20532}.

\bibitem{GoodToulmin}
I. J. Good and G. H. Toulmin.
The number of new species, and the increase in population coverage, when a
sample is increased.
\emph{Biometrika} \textbf{43}(1--2) (1956), 45--63.
\href{https://doi.org/10.1093/biomet/43.1-2.45}{doi:10.1093/biomet/43.1-2.45}.

\bibitem{Haeusler}
E. Haeusler.
On the rate of convergence in the central limit theorem for martingales with
discrete and continuous time.
\emph{Annals of Probability} \textbf{16}(1) (1988), 275--299.
\href{https://doi.org/10.1214/aop/1176991901}{doi:10.1214/aop/1176991901}.

\bibitem{Janson}
S. Janson.
On edge exchangeable random graphs.
\emph{Journal of Statistical Physics} \textbf{173} (2018), 448--484.
\href{https://doi.org/10.1007/s10955-017-1832-9}{doi:10.1007/s10955-017-1832-9}.

\bibitem{Mourrat}
J.-C. Mourrat.
On the rate of convergence in the martingale central limit theorem.
\emph{Bernoulli} \textbf{19}(2) (2013), 633--645.
\href{https://doi.org/10.3150/12-BEJ417}{doi:10.3150/12-BEJ417}.

\end{thebibliography}
\end{document}
